\documentclass[11pt]{article}
\usepackage[margin=1in]{geometry}
\usepackage[T1]{fontenc}
\usepackage[utf8]{inputenc}
\usepackage{lmodern}
\usepackage{microtype}
\usepackage{graphicx}
\usepackage{booktabs}
\usepackage{url}
\usepackage[hidelinks]{hyperref}
\usepackage{titlesec}
\usepackage{listings}
\usepackage{xcolor}

% Readable section spacing and paragraph layout
\titlespacing*{\section}{0pt}{10pt plus 3pt minus 2pt}{6pt}
\titlespacing*{\subsection}{0pt}{8pt plus 3pt minus 2pt}{4pt}
\setlength{\parskip}{6pt}
\setlength{\parindent}{0pt}
\linespread{1.03}

% Listings setup (more legible)
\lstset{
  basicstyle=\ttfamily\small,
  numbers=left,
  numberstyle=\tiny\color{gray},
  numbersep=8pt,
  frame=single,
  rulecolor=\color{black},
  breaklines=true,
  showstringspaces=false,
  columns=fullflexible,
  backgroundcolor=\color{black!1},
  xleftmargin=1.0em,
  xrightmargin=0.4em,
  aboveskip=8pt, belowskip=8pt
}

\title{CertTalk: A Clear, Proof-First Protocol for Two Agents With Private Maps}
\author{Agent Talk Team}
\date{October 27, 2025}

\begin{document}
\maketitle

\textbf{Abstract}\\
Two agents each see only part of a grid’s obstacles. They must decide if a route exists from start $s$ to goal $t$ in the union map $U$—and hand over a proof anyone can replay later. CertTalk does not ship maps; it ships a certificate. If a path exists, the certificate is the path. If no path exists, the certificate is a small set of blocked cells that separates $s$ from $t$. This writeup presents the idea in plain language, shows two worked examples, and prints full, annotated transcripts for the baselines so a reader can verify what each system actually sends.

\section{What We Decide (and What We Don’t)}
\textbf{Problem.} A knows obstacles $A$. B knows obstacles $B$. Neither sees the union $U = A\cup B$. The question is: does a 4-neighbour path exist from $s$ to $t$ in $U$?

\textbf{Contract.} The dialogue ends with a single certificate that a third-party checker (the “oracle”) can replay using only the certificate, $s$, $t$, grid size, and $U$. The oracle never reads the negotiation transcript and never needs the agents’ planner code.

\textbf{Out of scope.} No motion control, dynamics, or multi-agent coordination beyond the yes/no reachability proof. Certificates are short (target $\le$ 256 bytes) and archivable.

\section{How the Proof Works}
\textbf{Path certificate (checker’s steps).}
\begin{enumerate}
  \item Decode Base64 runs (RLE4 bytes).\label{step:path-decode}
  \item Verify \texttt{digest16} (CRC16 of raw bytes).\label{step:path-crc}
  \item Expand runs into coordinates from $s$ to $t$; check 4-connectivity and bounds.\label{step:path-expand}
  \item Confirm no coordinate lies in a blocked cell of $U$.\label{step:path-free}
  \item Accept. The bytes \emph{are} the proof.
\end{enumerate}

\textbf{Cut certificate (checker’s steps).}
\begin{enumerate}
  \item Decode Base64 cell list (DELTA16) and witness bits.\label{step:cut-decode}
  \item Verify \texttt{digest16} (CRC16 of raw bytes).\label{step:cut-crc}
  \item Check each listed cell is blocked in $U$.\label{step:cut-cells}
  \item Remove the cells from the underlying 4-neighbour graph; verify $s$ and $t$ are disconnected.\label{step:cut-disc}
  \item Accept. The bytes \emph{are} the proof.
\end{enumerate}

\vspace{6pt}
\textit{Witness bits.} One bit per cut cell says “who saw this cell” (0=A, 1=B). They guide the responder during negotiation. The oracle ignores witness semantics except to check length.

\section{Worked Example: Reachable (seed 247)}
When a path exists, A proposes a concrete route. B replays it locally. If every step is free for B, B upgrades it to a certificate. Three messages end the dialogue.

\textbf{CertTalk — PATH success}
\begin{lstlisting}
{
  "t": "PATH_PROPOSE",
  "s": "v1", "q": 0,
  "p": { "runs": "QoVBgkKCRA==", "digest16": 22270 },
  "c": 16269
}
{
  "t": "PATH_CERT",
  "s": "v1", "q": 0,
  "p": { "runs": "QoVBgkKCRA==", "digest16": 22270, "signed_by": ["B"] },
  "c": 49459
}
{
  "t": "ACK",
  "s": "v1", "q": 1,
  "p": { "ack_of": "PATH_CERT", "digest16": 22270 },
  "c": 3754
}
\end{lstlisting}
\textit{What B checks.} Decode runs; verify CRC; expand steps; confirm every coordinate is free in B’s map; certify. The certificate itself is the replayable proof.

\section{Worked Example: Unreachable (seed 145)}
When no path exists, A proposes a short list of blocked cells that separates $s$ from $t$. B validates per-cell (using witness attribution to allow “I saw it” from A) and then validates disconnection. Three messages suffice on this seed.

\textbf{CertTalk — CUT success}
\begin{lstlisting}
{
  "t": "CUT_PROPOSE",
  "s": "v1", "q": 0,
  "p": { "cs": "AAACAAH/Af8=", "k": 3, "digest16": 11880, "w": "AA==" },
  "c": 51694
}
{
  "t": "CUT_CERT",
  "s": "v1", "q": 0,
  "p": { "cs": "AAACAAH/Af8=", "k": 3, "digest16": 11880, "w": "AA==", "sb": ["B"] },
  "c": 43260
}
{
  "t": "ACK",
  "s": "v1", "q": 1,
  "p": { "ack_of": "CUT_CERT", "digest16": 11880 },
  "c": 27687
}
\end{lstlisting}
\textit{What B checks.} Decode cells and witness; verify CRC; accept each cell (blocked locally or credited to A); remove the set and confirm $s$ and $t$ are disconnected; certify.

\section{Baselines: What They Send (and Why It Matters)}
Different systems reach the same decision with different bytes. Below we show one clean unreachable seed (145) where all four succeed, and why a reader should care.

\textbf{Send-All — CUT success (seed 145)}
\begin{lstlisting}
{
  "type": "PROBE", "schema": "v1", "seq": 0,
  "payload": { "what": "CELLS",
    "indices": [2,3,4,11,13,14,17,19,20,22,27,32,36,37,38,39,46,47,63,65,68,70,71,78,80,83,85,95,98],
    "chunk": 0, "total": 1, "final": true },
  "crc16": 40037
}
{
  "type": "CUT_CERT", "schema": "v1", "seq": 0,
  "payload": { "cells": "AgAAAP4BAQA=", "encoding": "DELTA16_v1",
    "digest16": 34702, "k": 3, "witness": "Ag==", "signed_by": ["B"] },
  "crc16": 12256
}
{
  "type": "ACK", "schema": "v1", "seq": 1,
  "payload": { "ack_of": "CUT_CERT", "digest16": 34702 },
  "crc16": 57117
}
\end{lstlisting}
\textit{Why it matters.} Send-All may stream raw blocked indices first. The later certificate—not the raw indices—is the replayable proof.

\textbf{Greedy-Probe — CUT success (seed 145)}\quad Same as CertTalk above on this clean seed (three messages: CUT\_PROPOSE, CUT\_CERT, ACK). Probes appear only when needed.

\textbf{Responder-MinCut — CUT success (seed 145)}
\begin{lstlisting}
{
  "t": "SCHEMA", "s": "v1", "q": 0,
  "p": { "cut_enc": "DELTA16_v1", "path_enc": "RLE4_v1", "s": [0,0], "size": [10,10], "t": [9,9] },
  "c": 64959
}
{
  "t": "CUT_CERT", "s": "v1", "q": 0,
  "p": { "cs": "BwAHAAEAAQD9Af8B", "k": 5, "digest16": 4126, "w": "Hw==", "sb": ["B"] },
  "c": 32263
}
{
  "t": "NACK", "s": "v1", "q": 1,
  "p": { "n": "CUT_PROPOSE", "reason": "BLOCKED", "x": [[7,7],[9,7],[6,8]] },
  "c": 593
}
{
  "t": "PROBE", "s": "v1", "q": 1,
  "p": { "cs": [[7,7],[9,7],[6,8]], "wht": "CELLS" },
  "c": 55063
}
{
  "t": "HELLO", "s": "v1", "q": 2, "p": {}, "c": 26486
}
{
  "t": "DONE", "s": "v1", "q": 2, "p": {}, "c": 22322
}
\end{lstlisting}
\textit{Why it matters.} Shows the longer path to agreement when the first cut is insufficient; one small probe reconciles the disagreement.

\section{Results in One Table}
Medians across 1{,}000 cached instances (schema-free run).
\begin{center}
\begin{tabular}{lcc}
\toprule
System & Bytes & Rounds \\
\midrule
Send-All & 364 & 3 \\
CertTalk & 468 & 5 \\
Greedy-Probe & 468 & 5 \\
Responder-MinCut & 506 & 6 \\
\bottomrule
\end{tabular}
\end{center}
\noindent Takeaway: CertTalk spends about one extra hundred bytes versus Send-All, but those bytes \emph{are the proof} (path or cut), which simplifies auditing.

\section{Limits and What to Expect}
\begin{itemize}
  \item Two agents; $10\times 10$ grids; reliable in-order channel. Larger grids may require chunked certificates.
  \item No privacy model: witness bits reveal who attested each cut cell. Signatures are optional if authenticity is required.
  \item The checker never reads the transcript; it only replays the final artifact with $U$.
\end{itemize}

\end{document}
